% GENERATED FILE. Edit canonical lessons, then run npm run generate:books.
% canonical-source-hash: fnv1a64:fd170ba8e5a7ed86
% canonical-lessons: RU-C266-tuda-i-obratno, RU-C266-opozdal, RU-C266-review

\chapter{A Ticket, and a Missed Train}
\label{ch:ru-a-ticket-and-a-missed-train}
\hlchaptermodality{266}

\begin{chapteropening}
\textbf{By the end of this chapter:} \emph{I can buy a ticket, \ru{туда} \ru{и} \ru{обратно}, at the station, and report a problem: \ru{я} \ru{опоздал} \ru{на} \ru{поезд}, \ru{я} \ru{потерял} \ru{билет}.}

Complete RU-C266-review: buy a return ticket at the till, then report a missed train and a lost ticket.
\end{chapteropening}

\section[\texorpdfstring{\ru{билет} \ru{туда} \ru{и} … (bilét tudá i …)}{билет туда и … (bilét tudá i …)}]{\ru{билет} \ru{туда} \ru{и} \ru{обратно} (bilét tudá i obrátno) — a return ticket}

\label{lesson:RU-C266-tuda-i-obratno}



\begin{quote}
Say \textquotedblleft{}a ticket\textquotedblright{}, \textquotedblleft{}a train\textquotedblright{} and \textquotedblleft{}the till\textquotedblright{}. (\emph{Bilét. Póyezd. Kássa.}) Then ask the way to the till. (\emph{Kak proytí k kásse?})
\end{quote}

\subsection*{You'll want to know: \ru{туда} \ru{и} \ru{обратно}}
\textbf{\ru{туда} \ru{и} \ru{обратно}} — \emph{tudá i obrátno} — \textquotedblleft{}there and back\textquotedblright{}. At the till it is a return ticket:

\begin{quote}
\textbf{— \ru{Один} \ru{билет} \ru{до} \ru{Москвы}, \ru{пожалуйста}. — \ru{Туда} \ru{и} \ru{обратно}?} — \emph{odín bilét do Moskvý, pozhálusta. tudá i obrátno?} — One ticket to Moscow, please. Return? \textbf{— \ru{Да}, \ru{туда} \ru{и} \ru{обратно}.} — \emph{da, tudá i obrátno} — Yes, a return.
\end{quote}

A single is \textbf{\ru{в} \ru{одну} \ru{сторону}} — \emph{v odnú stóronu} — \textquotedblleft{}in one direction\textquotedblright{}.

\begin{grammarlens}[title={Grammar Lens: \ru{до} a place, \ru{на} a train}]
A ticket goes \textbf{\ru{до}} the place, as far as it, and the place changes as it did after \emph{\ru{до}} in the street: \emph{\ru{светофор}}, \emph{\ru{до} \ru{светофора}}; \emph{\ru{Москва}}, \textbf{\ru{до} \ru{Москвы}}.

A ticket is \textbf{\ru{на}} the train or the bus, and \emph{\ru{поезд}} and \emph{\ru{автобус}} keep their plain form: \textbf{\ru{билет} \ru{на} \ru{поезд}}, \textbf{\ru{билет} \ru{на} \ru{автобус}} — \emph{bilét na póyezd, bilét na avtóbus}.
\end{grammarlens}

\subsection*{The exchange}
\begin{quote}
\textbf{— \ru{Сколько} \ru{стоит} \ru{билет}? — \ru{Тысяча} \ru{рублей}.}
\end{quote}

\emph{Skólko stóit bilét?} — \textquotedblleft{}How much is the ticket?\textquotedblright{} \emph{Týsyacha rubléy} — \textquotedblleft{}A thousand roubles.\textquotedblright{} After \emph{\ru{тысяча}}, \emph{\ru{рубль}} takes the counting form \textbf{\ru{рублей}}.

\begin{cousinweb}
\textbf{\ru{обратно}} is the adverb of \emph{\ru{обратный}}, \textquotedblleft{}reverse, return\textquotedblright{}. EtymDB 2.1 lists \emph{\ru{обратный}} without a parent; by the usual account it is \emph{\ru{об}-}, \textquotedblleft{}around\textquotedblright{}, with the root of \textbf{\ru{вертеть}}, \textquotedblleft{}to turn\textquotedblright{}, from Proto-Indo-European \emph{*wert-}: a ticket that turns round. \textbf{\ru{туда}} is \textquotedblleft{}to that place\textquotedblright{}, on the same pointing stem as \emph{\ru{тот}}.
\end{cousinweb}

\subsection*{Your turn}
\begin{itemize}
\raggedright
  \item \emph{Say it:} \emph{Odín bilét do Moskvý, pozhálusta.}
  \item \emph{Say it:} \emph{Tudá i obrátno.} Then a single: \emph{V odnú stóronu.}
  \item \emph{Say it:} \emph{Dva biléta na póyezd.}
  \item \emph{Answer:} \emph{Skólko stóit bilét?} \emph{Týsyacha rubléy.}
\end{itemize}

\begin{tcolorbox}[breakable,colback=teal!4,colframe=teal!35!black,title={Before you move on}]
\textquotedblleft{}A return ticket\textquotedblright{}? (\emph{\textbf{\ru{Билет} \ru{туда} \ru{и} \ru{обратно}.}}) \textquotedblleft{}A ticket to Moscow\textquotedblright{}? (\emph{\textbf{\ru{Билет} \ru{до} \ru{Москвы}.}}) \textquotedblleft{}A ticket for the train\textquotedblright{}? (\emph{\textbf{\ru{Билет} \ru{на} \ru{поезд}.}})
\end{tcolorbox}
\section[\texorpdfstring{\ru{я} \ru{опоздал} \ru{на} \ru{поезд} (ya opozdál na …)}{я опоздал на поезд (ya opozdál na …)}]{\ru{я} \ru{опоздал} \ru{на} \ru{поезд} (ya opozdál na póyezd) — I've missed the train}

\label{lesson:RU-C266-opozdal}



\begin{quote}
Say \textquotedblleft{}to be late\textquotedblright{} and \textquotedblleft{}to lose\textquotedblright{}. (\emph{Opázdyvat. Teryát.}) Then \textquotedblleft{}a return ticket for the train\textquotedblright{}. (\emph{Bilét na póyezd tudá i obrátno.})
\end{quote}

\subsection*{You'll want to know: \ru{я} \ru{опоздал} \ru{на} \ru{поезд}}
\begin{quote}
\textbf{\ru{Я} \ru{опоздал} \ru{на} \ru{поезд}.} — \emph{ya opozdál na póyezd} — I've missed the train. A woman says \textbf{\ru{Я} \ru{опоздала} \ru{на} \ru{поезд}.}
\end{quote}

Russian says \textquotedblleft{}I was late \textbf{for} the train\textquotedblright{}, with \textbf{\ru{на}}, as on a ticket. \textbf{\ru{Опоздать}} is the perfective partner of \emph{\ru{опаздывать}}: one lateness, and the train is gone.

\begin{grammarlens}[title={Grammar Lens: a problem is a result}]
English often reports a problem with \textquotedblleft{}I have\textquotedblright{}: I've lost, it has gone. Russian uses the perfective past, because the result is still true now:

\begin{quote}
\textbf{\ru{К} \ru{сожалению}, \ru{я} \ru{потерял} \ru{билет}.} — \emph{k sozhaléniyu, ya poteryál bilét} — Unfortunately I've lost my ticket. \textbf{\ru{Поезд} \ru{уже} \ru{ушёл}.} — \emph{póyezd uzhé ushól} — The train has already left.
\end{quote}

\textbf{\ru{Потерял}} is the partner of \emph{\ru{терять}}; \textbf{\ru{ушёл}}, of \emph{\ru{уходить}}.
\end{grammarlens}

\subsection*{The exchange}
\begin{quote}
\textbf{— \ru{Извините}, \ru{у} \ru{меня} \ru{проблема}. \ru{Я} \ru{опоздал} \ru{на} \ru{поезд}.} \textbf{— \ru{Ничего}. \ru{Есть} \ru{поезд} \ru{через} \ru{час}.}
\end{quote}

\emph{U menyá probléma} — \textquotedblleft{}I have a problem\textquotedblright{}, literally \textquotedblleft{}at me, a problem\textquotedblright{}. \emph{Nichevó. Yest póyezd chérez chas} — \textquotedblleft{}Never mind. There's a train in an hour.\textquotedblright{} \textbf{\ru{Есть}}, \textquotedblleft{}there is\textquotedblright{}, is the present of \emph{\ru{быть}} that Russian does say aloud. When it is serious: \textbf{\ru{Помогите}, \ru{пожалуйста}!}

\begin{cousinweb}
\textbf{\ru{опоздать}} is \emph{\ru{о}-} with the root of \textbf{\ru{поздно}}, \textquotedblleft{}late\textquotedblright{}. EtymDB 2.1 lists \emph{\ru{опоздать}} without a parent; the usual account takes \emph{\ru{поздно}} back to a Proto-Slavic \emph{*pozdŭ} and relates it to Latin \emph{post}, \textquotedblleft{}after\textquotedblright{}.
\end{cousinweb}

\subsection*{Your turn}
Say these aloud:

\begin{itemize}
\raggedright
  \item \emph{Ya opozdál na póyezd}, then a woman: \emph{Ya opozdála na póyezd.}
  \item \emph{K sozhaléniyu, ya poteryál bilét.}
  \item \emph{Izvinítye, u menyá probléma. Póyezd uzhé ushól.}
  \item a problem of your own at a station, with one perfective verb
\end{itemize}

\begin{tcolorbox}[breakable,colback=teal!4,colframe=teal!35!black,title={Before you move on}]
\textquotedblleft{}I've missed the train\textquotedblright{}? (\emph{\textbf{\ru{Я} \ru{опоздал} \ru{на} \ru{поезд}.}}) \textquotedblleft{}I've lost my ticket\textquotedblright{}: which aspect? (\textbf{Perfective}, \emph{\textbf{\ru{потерял}}}: it is still lost.) \textquotedblleft{}I have a problem\textquotedblright{}? (\emph{\textbf{\ru{У} \ru{меня} \ru{проблема}.}})
\end{tcolorbox}
\section[(review)]{Review — a ticket, and a train that left without you}

\label{lesson:RU-C266-review}



\begin{quote}
Ask for one return ticket to Moscow. (\emph{Odín bilét do Moskvý, tudá i obrátno, pozhálusta.})
\end{quote}

\begin{grammarlens}[title={What you've built}]
Two lessons, and the station can be handled out loud. A ticket goes \textbf{\ru{до}} a place and \textbf{\ru{на}} a train; a return is \textbf{\ru{туда} \ru{и} \ru{обратно}}. And when it goes wrong, the perfective past says what happened and is still true: \textbf{\ru{я} \ru{опоздал} \ru{на} \ru{поезд}}, \textbf{\ru{я} \ru{потерял} \ru{билет}}.
\end{grammarlens}

\subsection*{The exchange}
\begin{quote}
\textbf{— \ru{Извините}, \ru{я} \ru{опоздала} \ru{на} \ru{поезд}. \ru{Как} \ru{пройти} \ru{к} \ru{кассе}?} \textbf{— \ru{Идите} \ru{прямо}, \ru{потом} \ru{налево}.}
\end{quote}

The whole trip so far in two lines: a problem in the perfective, a question about the way, and an answer in two steps.

\subsection*{Your turn}
\begin{itemize}
\raggedright
  \item \emph{Say it:} \emph{Izvinítye, ya opozdála na póyezd. Kak proytí k kásse?}
  \item \emph{Say it:} \emph{Dva biléta na póyezd, tudá i obrátno.}
  \item \emph{Say it:} \emph{K sozhaléniyu, ya poteryál bilét.}
  \item \emph{Answer:} \emph{Pochemú ty opozdál?} Explain what happened, then what followed, with \emph{poétomu}
  \item \emph{Answer:} the cashier asks \emph{Tudá i obrátno?} Answer for a single: \emph{Nyet, v odnú stóronu.}
\end{itemize}

\begin{tcolorbox}[breakable,colback=teal!4,colframe=teal!35!black,title={Before you move on}]
\textquotedblleft{}She's missed the train\textquotedblright{}? (\emph{\textbf{\ru{Она} \ru{опоздала} \ru{на} \ru{поезд}.}}) \textquotedblleft{}Two tickets for the bus\textquotedblright{}? (\emph{\textbf{\ru{Два} \ru{билета} \ru{на} \ru{автобус}.}})
\end{tcolorbox}
